% ECONOMICS_PROBLEM: {"id": "F12-267", "title": "Global-value-chain shock propagation", "jel_code": "F12", "source_id": "OP-0084"}
% !TeX program = xelatex
\documentclass[11pt,a4paper]{article}
\usepackage[margin=23mm]{geometry}
\usepackage{fontspec}
\setmainfont{Latin Modern Roman}
\usepackage{amsmath,amssymb,mathtools}
\usepackage{xcolor,enumitem,booktabs,longtable,array,xurl,hyperref}
\definecolor{muted}{HTML}{53636C}
\hypersetup{colorlinks=true,urlcolor=blue,linkcolor=blue}
\setlength{\parindent}{0pt}
\setlength{\parskip}{5pt}
\newcommand{\R}{\mathbb R}
\newcommand{\E}{\mathbb E}
\newcommand{\N}{\mathbb N}
\newcommand{\ind}{\mathbf 1}
\newcommand{\Var}{\operatorname{Var}}
\newcommand{\Cov}{\operatorname{Cov}}
\newcommand{\supp}{\operatorname{supp}}
\DeclareMathOperator*{\argmin}{arg\,min}
\DeclareMathOperator*{\argmax}{arg\,max}
\newcommand{\entrymeta}[3]{{\sffamily\small\color{muted}\textbf{\texttt{#1}}\quad|\quad #2\par #3\par}}
\newcommand{\Problemref}[1]{\texttt{#1}}
\newcommand{\JELref}[2]{\texttt{#2} (JEL \texttt{#1})}
\begin{document}
\section*{Global-value-chain shock propagation}
\textbf{Problem ID:} \texttt{F12-267}\quad \textbf{JEL:} \texttt{F12}\par
% BEGIN ECONOMICS STATEMENT
\label{problem:OP-0084}
\label{atlas:prob:T4}

\noindent\textbf{Atlas ID:} T4; \textbf{original number:} 84; \textbf{field:} International trade. \textbf{Classification:} Editorial research formulation (theoretical/empirical); not a certified unsolved theorem.

\subsection*{Definitions and assumptions}
Let firms $i=1,\ldots,n$ trade intermediate goods at dates $t=0,\ldots,H$. Production satisfies $y_{it}\le F_i(k_{it},\ell_{it},x_{1it},\ldots,x_{nit};z_{it})$, where $x_{jit}\ge0$ is input from supplier $j$, $z_{it}$ is productivity and $F_i$ is a specified monotone technology. Inventories obey feasible stock-flow accounting; establishing a supplier relationship costs resources and takes a stated time. Firms choose sourcing, production, storage and investment to maximize finite expected discounted profits. Household demand, goods clearing, network formation and a specified equilibrium selection jointly determine quantities and prices.

\subsection*{Formal research question}
For an exogenous disruption $z^1$ and a common baseline $z^0$, define aggregate value-added loss by
\[
 L_H(\theta,z^1)=\sum_{t=0}^{H}\delta^t
 [\mathrm{VA}_t(e_\theta(z^0))-\mathrm{VA}_t(e_\theta(z^1))],\quad 0<\delta\le1,
\]
using a common valuation convention and subtracting intermediate purchases to avoid double counting. Identify the technologies, inventory policies and supplier-switching costs needed to estimate the distribution of $L_H$ and its upper tail. Characterize conditions under which linear input-output propagation approximates the nonlinear equilibrium, and bound its approximation error as shock size, substitutability and available inventories vary.

\subsection*{Interpretation and scope}
Fixed input-output coefficients are a benchmark rather than a description of all production. Bottlenecks may imply nondifferentiable responses, so a local derivative need not represent large-disruption losses. Domestic losses and world losses differ when activity moves internationally. A causal disruption design must exclude simultaneously changing final demand or explicitly model it. Network observation should cover exits, missing transactions and supplier substitution; otherwise apparent recovery may be sample attrition. For tail prediction, specify a shock distribution and report sensitivity to poorly observed dependence across suppliers. Welfare requires utility and incidence beyond the value-added loss statistic. Inventory dynamics must include depreciation and delivery lags, and intermediate purchases should be recorded at consistently defined transfer prices before comparing aggregate paths.

\subsection*{Source audit and status}
[S1] and [S3], including linked source [43], concern production-sharing and value-added accounting; [S2] concerns organization and contracts. These sources do not establish a complete nonlinear shock-propagation solution with inventories and rewiring. This entry is an editorial extension, and the citations identify relevant foundations rather than certify a currently open theorem.

\subsection*{References}

\begin{enumerate}[label={[S\arabic*]},leftmargin=*]

\item Robert C. Johnson and Guillermo Noguera (2012). \emph{Accounting for Intermediates: Production Sharing and Trade in Value Added}. Journal of International Economics 86(2), 224–236. \url{https://www.sciencedirect.com/science/article/pii/S002219961100122X}.
\textit{Locator and role:} Publisher or author abstract / bibliographic record; Value-added rather than gross-trade accounting. Provides background or a benchmark result; does not establish that the editorial question remains unresolved in 2026.

\item Pol Antràs and Davin Chor (2013). \emph{Organizing the Global Value Chain}. Econometrica 81(6), 2127–2204. \url{https://doi.org/10.3982/ECTA10813}.
\textit{Locator and role:} Publisher or author abstract / bibliographic record; Contracting and organization along sequential production chains. Provides background or a benchmark result; does not establish that the editorial question remains unresolved in 2026.

\item Zhi Wang, Shang-Jin Wei, and Kunfu Zhu (2013). \emph{Quantifying International Production Sharing at the Bilateral and Sector Levels}. NBER Working Paper 19677. \url{https://doi.org/10.3386/w19677}.
\textit{Locator and role:} Publisher or author abstract / bibliographic record; Bilateral and sectoral production-sharing decompositions; resolves link [43]. Provides background or a benchmark result; does not establish that the editorial question remains unresolved in 2026.

\end{enumerate}

\subsection*{Common conventions: Source mathematical conventions}

Unless an entry states otherwise, agent and firm sets are finite. State and action spaces are measurable, strategies and outcome maps are measurable, probability laws are specified, and expectations exist for the targets under discussion. Infinite-horizon discount factors lie in $(0,1)$ and discounted objectives are finite. Norms, tolerances and complexity input sizes must be specified for computational comparisons. Constants or primitives that appear locally are inputs, not empirical facts asserted by this catalogue.

\textbf{Identification.} A statistical model $m\in\mathcal M$ implies an observable law $P_m$ and target $\theta(m)$. At law $P$, its identified set is
\[
\Theta(P;\mathcal M)=\{\theta(m):m\in\mathcal M,\ P_m=P\}.
\]
Point identification holds when this set is a singleton, or equivalently when $P_{m_1}=P_{m_2}$ implies $\theta(m_1)=\theta(m_2)$ for all compatible models. Sharp bounds mean bounds attained, or approached when the boundary is not attained, by compatible models. Writing this set is a definition, not a solution: the substantive task is to establish informative restrictions and derive or estimate the set. An unrestricted-confounding model may give only trivial bounds.

\textbf{Inference.} Identification, estimability and valid uncertainty quantification are separate questions. Fix $\alpha\in(0,1)$. A confidence procedure $C_n$ over a declared law class $\mathcal P$ has asymptotic uniform coverage at level $1-\alpha$ when
\[
\liminf_{n\to\infty}\inf_{P\in\mathcal P}\Pr_P\{\theta(P)\in C_n\}\geq1-\alpha.
\]
For set-identified targets the coverage event must be defined separately, for example coverage of the full identified set. If an entry uses identification-indexed models instead of uniquely indexed laws, the same coverage convention is applied over those models. Pointwise limits do not establish uniform validity, and asymptotic validity does not guarantee finite-sample precision. Sample sizes, dependence structures and weak-identification sequences are part of each application.

\textbf{Counterfactuals.} $Y_i(a)$ is a potential outcome under a specified intervention, including its timing, implementation and versions. It is not $Y_i$ conditional on voluntarily choosing $a$. A full intervention vector permits interference; shorthand scalar treatment notation does not silently impose no interference. Assignment, selection, attrition, anticipation and measurement must be specified. A claim about an instrument's local effect is not automatically a population policy effect.

\textbf{Economic equilibrium.} A counterfactual model includes best responses, constraints, feasibility, market clearing, state transitions and expectations consistent with the stated information. When several equilibria exist, either a declared selector is an input or the target is set-valued. Comparisons across policy regimes must use the stated selection convention. Reduced-form relationships are not presumed invariant after an equilibrium-changing intervention.

\textbf{Optimization and welfare.} Optimization is over the stated feasible set. If attainment is not established, $\sup$ or $\inf$ and $\varepsilon$-optimal choices replace an asserted maximizer or minimizer. For example, with value $V=\sup_{a\in\mathcal A}W(a)<\infty$, an $\varepsilon$-optimal set is $\{a\in\mathcal A:W(a)\geq V-\varepsilon\}$ for $\varepsilon>0$. Benefits, costs, risk and health or environmental outcomes can be aggregated only using declared units and conversion weights. Interpersonal utility weights, the treatment of fiscal transfers and foreign profits, discounting, and the behavioral welfare criterion are explicit inputs. A comparative-static derivative additionally requires existence and appropriate differentiability of the selected counterfactual.

\textbf{Sources.} Local labels [S1], [S2], and so forth restart in every question. Locators identify the relevant abstract, model, section or result at the level actually checked; a bibliographic or abstract-level verification is not represented as inspection of an entire proof. Working papers are distinguished from later publications and changing versions. Source summaries are paraphrased. All URL checks and editorial formulations refer to the audit date stated above.
% END ECONOMICS STATEMENT
\end{document}
